% Sources:
% - tamarin/rq-v2-minimal/rqv2_relaxed.spthy
% - tamarin/rq-v2-minimal/rqv2_message_dedup.spthy
% - tamarin/rq-v2-minimal/rqv2_party_admission.spthy
% - docs/paper/formal-model-draft.md
% - docs/rq-v2/g2-admission-semantics.md
% - docs/rq-v2/rq-v2-threat-model.md

\subsection{Modeling Objective and Abstraction}

Our model is not a symbolic verification of the complete \kwaay construction.
It isolates the identity relation applied when two sender-associated entries
are composed and admitted for one abstract \batchreceive execution. We choose
to omit cryptographic and application-layer detail in order to compare
admission predicates under explicit origin and freshness assumptions. This
choice does not establish that omitted computations cannot affect concrete
receiver behavior. The model retains explicit party,
sender-session, message, batch, and receiver-context coordinates while
abstracting the cryptographic computations that produce and process concrete
K-Waay messages.

The source interface associates each batch element with a party and groups
elements within one receiver invocation
\cite[Secs.~4.1 and~5.1]{collins2024kwaayfull}.
Table~\ref{tab:source-abstraction} separates these source-level relationships
from the assumptions used to encode them. The purpose of the reduction is to
keep party projection, repeated input, and batch membership visible while
varying the admission coordinate. It is not a derivation of the receiver's
cryptographic success conditions.

\begin{table}[htbp]
  \centering
  \caption{Source-to-abstraction mapping and its status.}
  \label{tab:source-abstraction}
  \small
  \begin{tabularx}{\linewidth}{@{}
    >{\raggedright\arraybackslash}X
    >{\raggedright\arraybackslash}X
    >{\raggedright\arraybackslash}X@{}}
    \toprule
    Source concept & Symbolic representation & Status \\
    \midrule
    Party associated with an input & Party coordinate $A$ &
      Source-derived semantic role; identity resolution is assumed \\
    Sender-associated batch input & Entry $(A,\mathit{sid},m)$ &
      Reduced tuple, not the wire format \\
    Elements in one receiver invocation & Shared
      $(\mathit{bid},\mathit{rst})$ & Model-local context identifiers \\
    Stated different-party condition & $A_1\neq A_2$ &
      Two-slot instance of the target relation \\
    Sender occurrence and message & Fresh $\mathit{sid}$ and $m$ &
      Prototype assumption, not concrete session-ID or message construction \\
    Component processing & \receiveraccept under an exact origin fact &
      Model-local observation, not a derived-key output \\
    \bottomrule
  \end{tabularx}
\end{table}

Public keys, prekey bundles, signatures, decapsulation, and batch-wide
cryptographic failure have no operational encoding here. The persistent
\texttt{!Sent} fact supplies an idealized origin relation rather than an
implementation of those checks. An admitted pair may fail to produce two
receiver events if an exact origin is absent; the abstraction does not
translate such a stalled path into a concrete failure result. We do not prove
a refinement theorem from full K-Waay executions to this abstraction, nor
that an abstract witness can be realized as a concrete K-Waay execution.
Our claims therefore concern the source-motivated admission relation under
the stated assumptions, not preservation of all protocol behavior.

We encode the abstraction in Tamarin, which represents executions by
multiset-rewriting rules and records trace-relevant action facts
\cite{meier2013tamarin}. The action facts let us state reachability,
correspondence, injectivity, and identity-distinction properties without
treating the model as executable K-Waay pseudocode. Rules model only the
lifecycle needed to expose the admission decision and its component-level
acceptance boundary.

The batch size is fixed at two. Since \distinctparty is a pairwise relation,
any violation in a larger batch contains at least one pair of positions with
equal party projections. Two slots therefore form the smallest batch
configuration in which the violating relation can be represented. This
justifies the witness shape, but it does not establish full-protocol behavior
or interactions among entries for arbitrary batch sizes.

We compare three controlled variants of one abstraction: relaxed admission,
a message-level restriction, and party-level admission. They share the same
sender-origin mechanism, symbolic network, two-slot collection, fresh batch
and receiver-context coordinates, and sequential component-processing path.
The principal analytical variation is the predicate applied after collection.
The guarded variants also contain their corresponding rejection branches and
inequality actions, so we do not claim that the theory files differ in only a
single line. Table~\ref{tab:admission-model-design} summarizes the design
without reporting any analysis outcome.

\begin{table}[htbp]
  \centering
  \caption{Controlled design of the batch-admission variants.}
  \label{tab:admission-model-design}
  \small
  \begin{tabularx}{\linewidth}{@{}
    >{\raggedright\arraybackslash}p{0.20\linewidth}
    >{\raggedright\arraybackslash}p{0.43\linewidth}
    >{\raggedright\arraybackslash}X@{}}
    \toprule
    Variant & Admission decision & Analytical role \\
    \midrule
    Relaxed admission & No inequality guard and no rejection branch &
      Invariant-removed baseline \\
    Message-level restriction & Admit if $m_1\neq m_2$; provide an
      equal-message rejection branch & Message-level comparison control \\
    Party-level admission & Admit if $A_1\neq A_2$; provide an equal-party
      rejection branch & Two-slot realization of \distinctparty \\
    \bottomrule
  \end{tabularx}
\end{table}

\subsection{Common Two-Slot Lifecycle}

Each entry is represented by the tuple
\begin{equation}
  E=(A,\mathit{sid},m),
  \label{eq:model-entry-tuple}
\end{equation}
where $A$ is a modeled protocol principal, $\mathit{sid}$ is a fresh
symbolic sender-session coordinate used to distinguish sender occurrences,
and $m$ is a fresh symbolic message coordinate in the abstraction. Party
creation records a persistent principal fact. A sender
occurrence then generates fresh $\mathit{sid}$ and $m$, emits the action
\(\mathsf{Send}(A,\mathit{sid},m)\), exposes
\(\langle A,\mathit{sid},m\rangle\) to the symbolic adversarial network, and
stores a persistent origin fact for exactly that tuple. Reuse of the party
fact permits one modeled principal to originate multiple tuples without
identifying the party with a session or message.

The symbolic adversary may later supply exposed tuples to the collection
rules. A fresh batch is scoped by a batch identifier $\mathit{bid}$ and a
receiver-context coordinate $\mathit{rst}$. The latter is a symbolic identity
for the shared receiver-side context; it is not a representation of the full
K-Waay receiver state. Collection advances through two linear states: the
first stores \(E_1\) and opens the second position, while the second stores
\(E_2\) and produces one collected pair under the same
\((\mathit{bid},\mathit{rst})\) context.

After collection, the selected variant either rejects the pair or creates the
state for processing the first component. An admitted path emits
\(\mathsf{BatchReceive}(\mathit{bid},\mathit{rst})\). Processing is sequential:
the first component can emit
\(\receiveraccept(A_1,\mathit{sid}_1,m_1,\mathit{bid},\mathit{rst})\) and then
opens the second component; the second can emit the analogous event for
\(E_2\) and reaches a completed-batch state. This is the symbolic lifecycle of
our abstraction, not a reconstruction of the complete K-Waay control flow.

Each processing step requires the persistent origin fact matching the full
tuple \((A_i,\mathit{sid}_i,m_i)\). Knowledge of the same party coordinate is
therefore insufficient to justify a receiver event for a different tuple.
Conversely, persistence permits the same exact origin fact to support more
than one processing step if the corresponding tuple is collected more than
once. The \receiveraccept action does not carry an explicit slot index; its
event occurrence and the linear processing state distinguish the two steps,
while $\mathit{bid}$ and $\mathit{rst}$ keep them in the same modeled context.

\subsection{Relaxed Admission Model}

The relaxed admission model is the invariant-removed baseline. Once both
entries have been collected, its admission rule emits
\(\mathsf{BatchReceive}(\mathit{bid},\mathit{rst})\) and starts component
processing without requiring either $A_1\neq A_2$ or $m_1\neq m_2$. It has no
post-collection rejection branch.

Accordingly, the admission rule itself does not rule out an exact repeated
tuple, nor does it rule out two distinct tuples that share a party coordinate.
This statement describes the rule's syntactic guard only. Whether either
composition participates in a complete trace satisfying a later analysis
property is deferred to Section~5.

\subsection{Message-Level Restriction Model}

The message-level restriction model changes the post-collection decision by
comparing $m_1$ and $m_2$. For a collected pair with syntactically equal
messages, the model provides a rejection branch that emits
\(\mathsf{Reject}(\mathit{bid},\mathit{rst})\) and reaches a modeled
batch-rejected state. The admitted branch records
\(\mathsf{Neq}(m_1,m_2)\) together with
\(\mathsf{BatchReceive}(\mathit{bid},\mathit{rst})\). A global restriction
excludes any trace containing \(\mathsf{Neq}(x,x)\), giving the action its
intended symbolic inequality meaning.

The admitted branch therefore requires $m_1\neq m_2$, but contains no
predicate requiring $A_1\neq A_2$. This variant is a controlled surrogate: it
constrains the message coordinate while leaving party equality unconstrained.
The logical non-implication
\[
  m_1\neq m_2
  \ \not\Longrightarrow\ 
  A_1\neq A_2
\]
motivates the comparison, but is not itself a model outcome. Nor is this
variant proposed as a concrete mitigation for K-Waay; it exists to test
whether a restriction over one identity coordinate can stand in for the
party-level objective.

This is an auxiliary control under the prototype's fresh-message abstraction,
not a general formalization of deployed deduplication. Among legitimate
sender-origin tuples, each fresh $m$ identifies one sender occurrence; message
equality therefore detects reuse of that origin in this encoding. Real-world
message encodings, collision behavior, normalization, and repeated sends with
identical content are not modeled. The control deliberately gives exact-message
comparison this favorable interpretation while leaving the party coordinate
separate.

\subsection{Party-Level Admission Model}

The party-level admission model applies the corresponding decision to
$A_1$ and $A_2$. For a collected pair with the same party coordinate, the
model provides a rejection branch that emits
\(\mathsf{Reject}(\mathit{bid},\mathit{rst})\) and reaches the rejected state.
The admitted branch records \(\mathsf{Neq}(A_1,A_2)\) and
\(\mathsf{BatchReceive}(\mathit{bid},\mathit{rst})\), subject to the same
global inequality restriction used by the message-level model. Message and
session coordinates are not substituted for the party comparison.

This rule is one executable realization of \distinctparty for the two-slot
abstraction: it directly constrains the coordinate named by the invariant.
It does not identify a unique enforcement mechanism for an implementation;
the same relation might instead be maintained by a caller, batch builder, or
another trusted component. Those deployment choices are outside the model.
Section~5 evaluates the properties and non-vacuity obligations produced by
this modeled realization rather than assuming them here.

\subsection{Events and Analysis Properties}

Five action-fact forms expose the common observation interface. The action
\(\mathsf{Send}(A,\mathit{sid},m)\) records one modeled sender origin.
\(\mathsf{BatchReceive}(\mathit{bid},\mathit{rst})\) records passage through
the admitted branch for one batch context; it is not the full protocol
implementation. A component-level
\(\receiveraccept(A,\mathit{sid},m,\mathit{bid},\mathit{rst})\) records
processing under an exact persistent sender-origin fact. The guarded variants
add \(\mathsf{Reject}(\mathit{bid},\mathit{rst})\) for their modeled
post-collection rejection decision and \(\mathsf{Neq}(x,y)\) for the
restricted inequality encoding. The completed and rejected batch markers are
state facts rather than action events.

Section~5 uses these observations for four property classes. First,
\emph{reachability} obligations ask whether representative admitted and
rejected behaviors, repeated-input compositions, and, where relevant, a valid
distinct-party batch have a trace. Such obligations are also needed to rule out
vacuous safety caused by admitting no batch. Second,
\emph{sender-origin correspondence} requires every \receiveraccept event to
have an earlier \(\mathsf{Send}\) event for the same complete
\((A,\mathit{sid},m)\) tuple.

Third, \emph{scoped occurrence injectivity} asks whether, in the presence of
a prior matching \(\mathsf{Send}\), two \receiveraccept occurrences agreeing
on \((A,\mathit{sid},m,\mathit{bid},\mathit{rst})\) must be the same event
occurrence. This property concerns one exact sender-origin tuple within one
batch and receiver context; it is not equivalent to same-batch party distinction.
Finally, \emph{identity-coordinate safety} asks whether different accepted
occurrences from one batch carry different messages in the message-restricted
variant or different parties in the party-restricted variant. These property
classes deliberately separate message identity, sender-occurrence
injectivity, and party identity. Their machine-checked outcomes, witness
traces, and comparative interpretation appear only in Section~5.
